\documentclass[10pt]{article}

% Pacotes extras necessários
\usepackage{amsmath}
\usepackage[lmargin=0.5in, rmargin=0.5in, tmargin=0.5in, bmargin=0.5in, includehead, includefoot]{geometry}
\usepackage{amsfonts}
\usepackage[utf8]{inputenc}
\usepackage[portuguese]{babel}
\usepackage{graphicx}
\usepackage{fancyhdr}
\usepackage{setspace}
\usepackage{listings}
\usepackage{url}
\usepackage{enumitem}
\usepackage{appendix}
\usepackage{subcaption}
\usepackage[dvipsnames]{xcolor}
\usepackage{tikz}
\usetikzlibrary{positioning}

% Define estilo para listagens de código
\lstdefinestyle{mystyle}{
    language=Matlab,
    inputencoding=latin1,
    extendedchars=true,
    backgroundcolor=\color{white},
    basicstyle=\ttfamily\footnotesize,
    numbers=left,
    numberstyle=\tiny\color{gray},
    numbersep=5pt,
    tabsize=2,
    breaklines=true,
    keywordstyle=\color{blue},
    commentstyle=\color{green!60!black},
    stringstyle=\color{orange},
    frame=single,
    keepspaces=true,
    showspaces=false,
    showstringspaces=false,
    showtabs=false,
}
\lstset{style=mystyle}

\graphicspath{{./images/}}

\setlength{\parindent}{0pt}
\setstretch{1.5}

\DeclareMathOperator{\sen}{sen}
\DeclareMathOperator{\sinc}{sinc}
\newcommand{\Lap}[1]{\mathcal{L}\left\{#1\right\}}
\newcommand{\bm}[1]{\boldsymbol{#1}}

% Cabeçalho e rodapé em todas as páginas
\pagestyle{fancy}
\fancyhf{}
\fancyhead[L]{Relatório das simulações}
\fancyfoot[C]{\thepage}

% Dados do Grupo
\title{
    Trabalho Nº - Gradiente Normalizado e Least-square Normalizado \\
    \large COE603 - Controle Adaptativo
}
\author{
    Caio Cesar Leal Verissimo - 119046624 \\
    Leonardo Soares da Costa Tanaka - 121067652 \\
    Lincoln Rodrigues Proença - 121076407 \\
    Engenharia de Controle e Automação - UFRJ \\
    Rio de Janeiro, Rio de Janeiro, Brasil \\
    Junho de 2025
}
\date{}

\begin{document}

\maketitle

% Geração automática do sumário
\tableofcontents
\newpage

\section{Descrição dos Algoritmos}

\subsection{Gradiente Normalizado}

\subsection{Least-square Normalizado}

\section{Resultados das Simulações}

\newpage

\section{Discussão dos Resultados}

Os resultados obtidos ao longo das simulações demonstram o bom desempenho dos algoritmos de Gradiente Normalizado (NG) e Least-Square Normalizado (NLS) em cenários variados, apesar de limitações inerentes ao tempo finito de simulação e à estrutura dos sinais de entrada.

No \textbf{caso nominal}, observou-se que o erro converge para valores próximos de zero em todas as parametrizações, mantendo-se limitado ao longo do tempo. Contudo, a solução ótima de identificação não foi plenamente atingida, refletindo-se nos picos de erro que ocorrem nas transições da onda quadrada. Tais descontinuidades impõem desafios à adaptação, uma vez que o sistema precisa reagir rapidamente a mudanças bruscas na referência.

Com a \textbf{mudança do sinal de referência} para uma forma com componente senoidal, notou-se uma leve melhoria na resposta adaptativa, atribuída à maior riqueza espectral do sinal. Entretanto, essa melhoria não foi suficiente para garantir a convergência aos parâmetros ótimos. Além disso, o erro passou a apresentar um comportamento mais ruidoso, o que é esperado, já que fora dos ganhos ideais a planta tem maior dificuldade em acompanhar sinais com maior conteúdo em frequência.

Ao considerar o \textbf{aumento do ganho de adaptação}, identificou-se uma melhora significativa na rapidez de convergência e uma redução no valor absoluto do erro. No entanto, esse ganho mais elevado pode comprometer a robustez do sistema, especialmente quando combinado a outras condições, como estados iniciais não nulos ou perturbações, podendo levar a respostas instáveis ou imprevisíveis.

Por fim, a influência das \textbf{condições iniciais} manifestou-se principalmente como um pico inicial no erro, que, após um período de transiente, tende a decair e a seguir o comportamento do caso nominal. Esse fenômeno destaca a importância de uma análise criteriosa na escolha dos ganhos de adaptação, já que a atualização dos parâmetros identificados é diretamente influenciada pelo valor do erro, especialmente nos instantes iniciais de simulação.

De forma geral, os experimentos reforçam a sensibilidade dos algoritmos NG e NLS a características como a forma do sinal de referência, os ganhos de adaptação e as condições iniciais do sistema. Apesar das limitações observadas, ambos os métodos mostraram-se eficazes em promover a adaptação dos parâmetros em tempo real, mantendo o erro de seguimento sob controle.

\end{document}
